A more general version of the next result has been announced \cite{JFR3}. We are very grateful to them for outlining their proof to us, which has inspired the argument that we give below.

\begin{Theorem}\label{thm:extensiontheory}
Let $G$ be a finite group, and $\mathfrak{C}$ be a fusion $2$-category. Then, faithfully $G$-graded extensions of $\mathfrak{C}$ are classified by homotopy classes of maps $\mathrm{B}G\rightarrow \mathrm{B}\mathscr{B}{\rm r}\mathscr{P}{\rm ic}(\mathfrak{C})$.
\end{Theorem}
\begin{proof}
For the purpose of this proof, it will be convenient to think of maps of spaces $\mathrm{B}G\rightarrow \mathrm{B}\mathscr{B}{\rm r}\mathscr{P}{\rm ic}(\mathfrak{C})$ as monoidal maps $G\rightarrow \mathscr{B}{\rm r}\mathscr{P}{\rm ic}(\mathfrak{C})$ of $(\infty,1)$-categories. We will also consider the $(\infty,1)$-category $\mathscr{C}$ obtained from $\mathbf{\operatorname{End}}_{\mathbf{F2C}}(\mathfrak{C})$ by only considering the invertible $n$-morphisms when $n\geq 2$. Then, by definition we have $\mathscr{B}{\rm r}\mathscr{P}{\rm ic}(\mathfrak{C})=\mathscr{C}^{\times}$ as monoidal $(\infty,1)$-categories.

We begin the proof by some general nonsense. Recall that algebras in the monoidal $(\infty,1)$-category $\mathscr{C}$ correspond precisely to lax monoidal functors $*\rightarrow \mathscr{C}$ (see \cite[Example~2.2.6.10]{L2}). We want a similar description for the notion of a (faithfully) $G$-graded algebra. In order to do so, consider the monoidal 1-category $G^{\sqcup}$, which is the coproduct completion of $G$. Now, there is a canonical algebra $A[G]$ in $G^{\sqcup}$ given by
 \begin{gather*}
A[G]:=\coprod_{g\in G}g,
\end{gather*}
 or equivalently a lax monoidal functor $*\rightarrow G^{\sqcup}$. Then, giving a $G$-grading on an algebra $A$ in $\mathscr{C}$ is equivalent to providing a factorization of lax monoidal functors
\begin{displaymath}
\begin{tikzcd}
* \arrow[rr,"A"] \arrow[rd, "A\lbrack G \rbrack"'] & & \mathscr{C} \\
 & G^{\sqcup}. \arrow[ru] &
\end{tikzcd}\end{displaymath}
But, given that $\mathscr{C}$ has direct sums, lax monoidal functors $G^{\sqcup}\rightarrow \mathscr{C}$ correspond exactly to lax monoidal functors $G\rightarrow \mathscr{C}$.

Now, given a map of spaces $\mathrm{B}G\rightarrow \mathrm{B}\mathscr{B}{\rm r}\mathscr{P}{\rm ic}(\mathfrak{C})$, or equivalently a strongly monoidal functor $F\colon G\rightarrow \mathscr{C}$, we obtain a $G$-graded algebra $\mathfrak{D}:=F(A[G])$ in $\mathfrak{C}$. We also write
\begin{gather*}
\mathfrak{D}=\boxplus_{g\in G}\mathfrak{D}_g=\boxplus_{g\in G}F(g).
\end{gather*}
 In fact, as $F(e) = \mathfrak{C}$ by definition, the monoidal 2-category $\mathfrak{D}$ is a faithfully $G$-graded extension of $\mathfrak{C}$. It is therefore enough to prove that $\mathfrak{D}$ is a fusion 2-category. The only property that is not obvious is rigidity. We have that $\mathbf{\operatorname{End}}_{\mathfrak{C}}(\mathfrak{D})$ is a fusion 2-category thanks to Theorem~\ref{thm:Moritadual} above. Let us fix $g\in G$. For any simple object $X$ in $\mathfrak{D}_g$, we can consider the left $\mathfrak{C}$-module 2-functor~${R_X\colon \mathfrak{D}\rightarrow \mathfrak{D}}$ given by $D\mapsto D\Box X$ on ${\mathfrak{D}_e\simeq \mathfrak{C}}$, and zero on $\mathfrak{D}_h$ for any~${e\neq h\in G}$. As the objects of the monoidal 2-category $\mathbf{\operatorname{End}}_{\mathfrak{C}}(\mathfrak{D})$ have duals, we can consider the right adjoint~$R_X^*$ of $R_X$. But, we have that $\mathfrak{D}_{g^{-1}}\rightarrow \mathbf{Fun}_{\mathfrak{D}_e}(\mathfrak{D}_g,\mathfrak{D}_e)$ is an equivalence of $\mathfrak{D}_e$-$\mathfrak{D}_e$-bimodule 2-categories as~$F$ is strongly monoidal, and therefore preserves duals. Thus, we find that ${R_X^*\simeq Y\Box (-)\colon \mathfrak{D}_g\rightarrow \mathfrak{D}_e=\mathfrak{C}}$ for some $Y$ in $\mathfrak{D}_{g^{-1}}$. This proves that $X$ has a left dual $Y$ in~$\mathfrak{D}$. One shows analogously that $X$ has a right dual.

Conversely, given $\mathfrak{D}=\boxplus_{g\in G}\mathfrak{D}_g$ a $G$-graded extension of $\mathfrak{C}= \mathfrak{D}_e$, we can consider the corresponding lax monoidal functor $F\colon G\rightarrow \mathscr{C}$. More precisely, we set $F(g) = \mathfrak{D}_g$ with lax monoidal structure given by the monoidal structure of $\mathfrak{D}$. We want to show that $F$ is strongly monoidal and factors through $\mathscr{B}{\rm r}\mathscr{P}{\rm ic}(\mathfrak{C})\subset \mathscr{C}$. It follows from the definition of an extension that $F$ is strongly unital. Proposition \ref{prop:gradedMoritaequivalence} above establishes that $F$ has the remaining desired properties.
\end{proof}
